% !TEX root = ../main.tex
\chapter*{General Introduction}
\phantomsection
\addcontentsline{toc}{chapter}{General Introduction}
\markboth{General Introduction}{General Introduction}
\label{chap:introduction}

Large Language Models (LLMs) have demonstrated strong capabilities across a
wide range of natural-language tasks. However, their answers remain constrained
by the knowledge encoded in their parameters and by the information available
in their input context. Retrieval-Augmented Generation (RAG) addresses this
limitation by complementing the language model with information retrieved from
an external corpus. This paradigm is particularly relevant to
knowledge-intensive question answering, where producing a correct answer
depends not only on the reasoning capabilities of the model, but also on whether
the necessary evidence is successfully retrieved and made available to it.

This dependence becomes more challenging in multi-hop question answering.
Unlike a direct factual question, a multi-hop question requires several
dependent pieces of information to be identified and combined. Some supporting
evidence may have only a weak lexical or semantic relation to the original
question, while information recovered at one reasoning step may be needed to
formulate the next information need. As a result, a retrieval system operating
only on the original question may fail to expose all of the evidence required
for the complete reasoning chain.

Multi-hop questions also differ substantially in their compositional demands.
Some can be resolved through relatively limited reasoning, whereas others
require several dependent operations and additional retrieval steps. Applying
the same processing strategy to every input may therefore be inefficient:
simple questions may receive unnecessary computation, while more compositional
questions may require additional reasoning or retrieval mechanisms. This
motivates adaptive approaches in which the processing strategy can depend on
the demands of the question.

One possible mechanism for handling compositional questions is question
decomposition. Instead of processing the original question as a single
information need, decomposition transforms it into a sequence of simpler
subquestions whose intermediate answers support later steps. These focused
subquestions can expose information needs that were implicit in the original
formulation and can therefore provide more targeted retrieval queries.
Nevertheless, generating a plausible decomposition does not guarantee that it
is useful for retrieval. A decomposition may omit an information need, retrieve
the required evidence only at large depths, or propagate an incorrect
intermediate answer to subsequent subquestions. Final-answer accuracy alone
does not reveal which of these events occurred.

This thesis studies adaptive multi-hop question answering from two
complementary perspectives. First, it investigates whether the compositional
difficulty of a question is reflected in the internal representations of an LLM
before answer generation, providing a potential signal about the amount of
processing the question may require. Then, once decomposition is selected as a
strategy for handling a compositional question, the thesis focuses on the
resulting retrieval process: how the generated subquestions can be evaluated
according to the evidence they retrieve, how decomposition quality can be
improved, and how a limited retrieval budget can be distributed more
effectively across the subquestions.

Following this progression, the work presented in this thesis makes three main
contributions:

\begin{enumerate}

    \item \textbf{Analysis of compositional difficulty in LLM
    representations.}
    The thesis investigates whether benchmark-defined compositional difficulty
    is linearly decodable from the hidden representations produced while a
    language model processes a question. The analysis is conducted across
    several language models and multi-hop question-answering benchmarks, and is
    complemented by controls involving label permutation, surface features,
    ordinal structure, and cross-dataset transfer.

    \item \textbf{Evidence-aware evaluation of question decomposition.}
    The thesis develops a retrieval-oriented framework for evaluating generated
    decompositions independently of final-answer correctness. The framework
    measures whether the annotated supporting evidence is retrieved, how deeply
    it appears in the resulting rankings, and how supporting passages can be
    associated with the subquestions responsible for retrieving them. This
    enables retrieval failures, decomposition behaviour, and answer-generation
    failures to be analysed separately.

    \item \textbf{Optimisation of decomposition and retrieval allocation.}
    The retrieval-oriented measurements are then used to study automated search
    over decomposition instructions and the allocation of a fixed passage
    budget across generated subquestions. The thesis further develops a
    structured evidence-feedback representation that organises retrieval,
    ownership, and intermediate-answer information for reflective prompt
    optimisation.
    
\end{enumerate}

Together, these contributions approach multi-hop question answering from both
the question and retrieval sides. The representation study examines whether
the model exposes information about compositional demand before answering,
while the decomposition studies examine how additional retrieval and reasoning
can be organised once decomposition is used. The resulting perspective treats
adaptation not only as the choice of whether additional processing may be
needed, but also as the problem of ensuring that this processing exposes the
required evidence efficiently.

The remainder of this report is organised as follows. Chapter~1 presents the
institutional and research context of the internship, together with the problem
statement, objectives, expected deliverables, and working methodology.
Chapter~2 reviews the background and state of the art in
retrieval-augmented generation, multi-hop question answering, retrieval
methods, adaptive computation, probing of internal representations, question
decomposition, and automated prompt optimisation.

Chapter~3 investigates whether compositional difficulty can be identified from
LLM representations before answer generation. Building on the need for
additional processing in compositional questions, Chapter~4 considers the case
in which question decomposition is used and develops an evidence-aware
framework for evaluating the retrieval behaviour of the resulting subquestions.
Chapter~5 then investigates how this behaviour can be improved through
automated decomposition-instruction search and learned passage allocation, and
presents a structured evidence-feedback interface for subsequent reflective
optimisation. Finally, the general conclusion summarises the main findings,
discusses the limitations of the work, and outlines directions for future
research.